\SourceLeaf{456}%
\noindent die nach Einsetzen der mit $\mathfrak{E},\mathfrak{F},\mathfrak{G}$ und $\mathfrak{E}',\mathfrak{F}',\mathfrak{G}'$ bezeichneten Werte zu folgender Gleichung führen:
\[
\left(\frac{\partial u}{\partial\zeta'}\right)^2=\left(\frac{\partial u'}{\partial\zeta}\right)^2.
\]

\noindent Demnach stellt der lineare Differentialausdruck
\[
ud\zeta\pm u'd\zeta'=d\varphi
\]
\noindent das totale Differential einer Funktion $\varphi$ der Variablen $\zeta,\zeta'$ dar. Diese Eigenschaft kann jedoch für je zwei aufeinander abwickelbare Flächen $S$ und $S'$ nur gelten, wenn das untere Vorzeichen gewählt wird. Bei entgegengesetzter Annahme müsste nämlich, wenn $S$ mit $S'$ zusammenfällt, das Produkt $2ud\zeta$ für jede Fläche ein totales Differential sein. Das tritt nur bei abwickelbaren Flächen ein.

Die vorstehende Gleichung enthält den von \textsc{Gauss} in Nr.\ 2 gewählten Ausgangspunkt für die allgemeine Lösung des Deformationsproblems.

Diese Gleichung gilt für je zwei aufeinander abwickelbare Flächen $S$ und $S'$, aber \textls[140]{nicht nur} für solche.

Wie \textsc{Gauss} sie für diese allgemeine Lösung nutzen wollte, lässt sich aus Nr.\ 4 der Notiz nicht unmittelbar erkennen. Diese Frage erfordert noch eingehende Untersuchungen der Eigenschaften der im jeweiligen Fall zu wählenden Funktion $\varphi$.
\begin{flushright}\textsc{Weingarten}.\end{flushright}

